\chapter{Sector torsor incidence: канонический \texorpdfstring{\(K_{2,2}\)}{K2,2} и blind no-go}

Четыре observed Yukawa sectors имеют естественные бинарные labels
\begin{equation}
 (q,h)\in\mathbb F_2^2,
\end{equation}
где \(q=0,1\) различает lepton/quark, а \(h=0,1\) --- coupling к
\(H/\widetilde H\):
\begin{equation}
 e=(0,0),\quad\nu=(0,1),\quad d=(1,0),\quad u=(1,1).
\end{equation}

\section{Новая идентификация четырёх minima}

Четыре transposition operators, ранее оставшиеся после family-square
selector, меняют quark/lepton bit и соединяют пары
\begin{equation}
 \nu d,\qquad\nu u,\qquad ed,\qquad eu.
\end{equation}
Следовательно, они являются ровно четырьмя edges полного bipartite graph
\begin{equation}
 K_{2,2}=\{e,\nu\}\star\{d,u\}.
\end{equation}

\begin{theorem}[Sector-torsor identification]
Четырёхэлементная transposition orbit не является произвольным menu из
четырёх family directions. После representation readout она канонически
совпадает с edge set quark--lepton sector torsor \(K_{2,2}\).
\end{theorem}

Это усиливает связь family carrier с observed representation, но edge
operators ещё требуется превратить в четыре vertex operators.

\section{Ориентированный incidence menu}

Пусть \(E_{\ell q}\) --- operator edge между lepton bit \(\ell\) и quark bit
\(q\). Для sign cocycle \(\epsilon_{\ell q}=\pm1\) degree-normalized vertex
incidence задаётся формулами
\begin{align}
 H_{\ell}&=-\frac12\sum_q\epsilon_{\ell q}E_{\ell q},\\
 H_q&=\frac12\sum_\ell\epsilon_{\ell q}E_{\ell q}.
\end{align}
Знак перед lepton/quark vertices есть обычная orientation incidence, а
\(1/2\) фиксирован общей степенью vertices. Полное discrete menu содержит
\(2^4=16\) cocycles и не имеет continuous sector coefficient.

Для каждого cocycle допускается один прежний общий ratio \(t\) в
\(R_4-tR_2\). Cocycles и \(t\) выбираются только по четырём quark mass ratios;
CKM остаётся blind.

\section{Mass train}

Лучший candidate имеет
\begin{equation}
 (\epsilon_{00},\epsilon_{01},\epsilon_{10},\epsilon_{11})
 =(-1,-1,+1,+1),
 \qquad
 t=19.944681.
\end{equation}
Его quark patterns равны
\begin{align}
 u&=(1.47982\times10^{-5},\,0.186971,\,1),\\
 d&=(0.000567355,\,0.179519,\,1),
\end{align}
с multiplicative errors
\begin{equation}
 (1.18,\,25.40,\,1.97,\,8.03).
\end{equation}
Как и в one-ratio gate, первая генерация может стать лёгкой, но обе middle
generations остаются слишком тяжёлыми.

\section{CKM blind}

После заморозки candidate получаем
\begin{equation}
 |V_q|\simeq
 \begin{pmatrix}
 0.96068&0.25521&0.10940\\
 0.25773&0.67639&0.68998\\
 0.10334&0.69091&0.71551
 \end{pmatrix}.
\end{equation}
Первый angle имеет правильный порядок:
\begin{equation}
 \frac{s_{12}}{s_{12}^{\rm CKM}}=1.14.
\end{equation}
Но два остальных ratios равны \(16.59\) и \(29.31\), а
\begin{equation}
 \frac{|J_q|}{J_{\rm CKM}}=58.39.
\end{equation}

\begin{theorem}[Oriented sector-incidence no-go]
Канонический sector torsor и исчерпывающее \(16\)-элементное orientation menu
не закрывают flavour sector даже с одним общим trained ratio. Mass-selected
branch проваливает middle-generation hierarchy, \(s_{23}\), \(s_{13}\) и
Jarlskog invariant.
\end{theorem}

\begin{remark}
Близкий Cabibbo angle является структурной зацепкой, но не pass. Он возникает
в том же ledger, где два других angles и CP invariant ошибочны на один--два
порядка.
\end{remark}

\begin{nogocriterion}[Запрет CKM-selected cocycle]
Нельзя перейти к другому из шестнадцати cocycles после открытия CKM. Menu был
заморожен и выбран по mass train; его blind failure закрывает всю объявленную
oriented-incidence ветвь.
\end{nogocriterion}

Следующим остаётся третий архитектурно отличный вариант: общий messenger
Schur complement, где sector-relative matrices возникают как коррелированные
products одних и тех же edge blocks, а не как vertex sums. Машинный ledger
сохранён в \path{s2t_v4_sector_torsor_incidence_gate_results.json}.